\documentclass[12pt,a4paper]{article}

\usepackage[romanian]{babel}
\usepackage[T1]{fontenc}
\usepackage[utf8]{inputenc}
\usepackage{lmodern}
\usepackage{enumitem}

\title{Indexarea în SGBD \\ Tipuri de Indexuri și Access Paths}
\author{}
\date{}

\begin{document}

\maketitle

\section{Indexarea în SGBD}
Indexarea este un mecanism de accelerare a operațiilor de căutare, sortare și 
filtrare în tabele. Un index funcționează similar unui tabel auxiliar, organizat 
într-o structură optimizată pentru acces rapid.

\section{Tipuri de Indexuri}

\subsection*{Index Clustered}
\begin{itemize}[noitemsep]
    \item determină ordinea fizică a rândurilor din tabel;
    \item există \textbf{un singur} index clustered per tabel;
    \item nu necesită memorie suplimentară pentru stocarea rândurilor, deoarece 
          datele sunt ordonate direct în tabel;
    \item structura tipică: \textbf{B-Tree} sau \textbf{B+Tree}.
\end{itemize}

\subsection*{Index Nonclustered}
\begin{itemize}[noitemsep]
    \item nu afectează ordinea fizică a datelor;
    \item stochează separat o structură de tip \textbf{B+Tree};
    \item necesită \textbf{memorie suplimentară} pentru index entries;
    \item fiecare nod frunză conține: \textit{(cheie, pointer către rând)}.
\end{itemize}

\section{Structuri de Indexare}

\subsection*{B-Tree}
\begin{itemize}[noitemsep]
    \item arbore echilibrat;
    \item nodurile interne conțin chei și pointeri;
    \item frunzele pot conține date sau pointeri către date.
\end{itemize}

\subsection*{B+Tree}
\begin{itemize}[noitemsep]
    \item toate datele se află în nodurile frunză;
    \item nodurile interne conțin doar chei pentru navigare;
    \item frunzele sunt legate secvențial, optim pentru \textbf{range scans}.
\end{itemize}

\section{Index Entry}
Un \textbf{index entry} reprezintă o înregistrare dintr-un nod frunză al indexului.

Conține:
\begin{itemize}[noitemsep]
    \item valoarea cheii indexate;
    \item pointer către rând (pentru nonclustered);
    \item pointer către pagina de date (pentru clustered).
\end{itemize}

\section{Access Paths}
Access paths reprezintă modurile în care un SGBD poate accesa datele folosind 
indexuri sau scanări directe.

\subsection*{Index Unique Scan}
\begin{itemize}[noitemsep]
    \item folosit când indexul este \textbf{unic};
    \item returnează cel mult un rând;
    \item cost minim, acces direct în arbore.
\end{itemize}

\subsection*{Index Range Scan}
\begin{itemize}[noitemsep]
    \item folosit pentru intervale: \texttt{WHERE x BETWEEN a AND b};
    \item navigare în B+Tree până la primul nod relevant, apoi parcurgere secvențială;
    \item eficient pentru date ordonate.
\end{itemize}

\subsection*{Index Full Scan}
\begin{itemize}[noitemsep]
    \item parcurge toate nodurile frunză ale indexului;
    \item util când indexul acoperă toate coloanele necesare (covering index);
    \item mai rapid decât un full table scan dacă indexul este mai mic decât tabelul.
\end{itemize}

\subsection*{Table Full Scan}
\begin{itemize}[noitemsep]
    \item parcurge întregul tabel;
    \item folosit când nu există index util sau filtrarea este slab selectivă.
\end{itemize}

\section{Avantaje și Dezavantaje ale Indexării}

\subsection*{Avantaje}
\begin{itemize}[noitemsep]
    \item accelerarea interogărilor;
    \item optim pentru căutări, sortări și join-uri;
    \item reduce costul operațiilor de acces.
\end{itemize}

\subsection*{Dezavantaje}
\begin{itemize}[noitemsep]
    \item consum de memorie suplimentară (nonclustered);
    \item încetinirea operațiilor \texttt{INSERT}, \texttt{UPDATE}, \texttt{DELETE};
    \item necesitatea întreținerii structurii arborelui.
\end{itemize}

\end{document}
